\documentclass[11pt,a4paper]{article}
\usepackage[T1]{fontenc}
\usepackage{lmodern}
\usepackage[margin=23mm]{geometry}
\usepackage{amsmath,amssymb,amsthm,mathtools,booktabs,array}
\usepackage[dvipsnames]{xcolor}
\usepackage[colorlinks=true,linkcolor=MidnightBlue,urlcolor=MidnightBlue,citecolor=MidnightBlue]{hyperref}
\usepackage{microtype,fancyhdr}
\hypersetup{pdftitle={Erdos Problem 708: Synthesis and Further Partial Results},pdfauthor={Research working note prepared for Denis Paliy},pdfsubject={Partial results; unrestricted exact and asymptotic conjectures remain unresolved}}
\pagestyle{fancy}
\fancyhf{}
\fancyhead[L]{\small Erd\H{o}s 708: synthesis and further partial results}
\fancyhead[R]{\small 23 November 2026}
\fancyfoot[C]{\thepage}
\setlength{\headheight}{14pt}
\setlength{\parskip}{4pt}
\setlength{\parindent}{0pt}
\newtheorem{theorem}{Theorem}[section]
\newtheorem{lemma}[theorem]{Lemma}
\newtheorem{corollary}[theorem]{Corollary}
\newtheorem{proposition}[theorem]{Proposition}
\theoremstyle{remark}
\newtheorem{remark}[theorem]{Remark}
\DeclareMathOperator{\lcm}{lcm}
\title{\textbf{Erd\H{o}s Problem 708:\\Synthesis and Further Partial Results}\\[5pt]
\large Prime-bank reductions, diffuse weights,\\dense smooth supports, and separated prime bands}
\author{Research working note prepared for Denis Paliy\\
\small AI-assisted derivations with separate internal mathematical audits}
\date{23 November 2026}
\begin{document}
\maketitle
\begin{center}
\fcolorbox{MidnightBlue}{blue!3}{\parbox{0.92\linewidth}{
\textbf{Status.} Neither the unrestricted exact bound \(2n\) nor the
unrestricted asymptotic bound \((2+o(1))n\) is proved here. The principal
new covering result treats all source prime factors below
\(\exp(o(\sqrt{\log M}))\). A separate exact theorem allows one large
prime per source together with a growing set of much smaller primes.
The proofs received internal cross-checks; external review, formal
verification, and priority are not claimed.
}}
\end{center}
\begin{abstract}
We consolidate the submitted \(32n^{5/2}\) argument with an earlier
cofactor-matching approach, and distinguish both from the already
available quadratic theorem of Ziming Wang. We prove structural
restrictions on a minimal counterexample by using forced prime-deletion
banks and total prime demands. On the analytic side, a partial
prime-power crossing argument and the fundamental lemma of the sieve
give a uniform approximate hinge inequality for arbitrarily many
sufficiently small raw weights. A controlled exceptional block incurs
an explicit defect. A sieve approximation of exact smooth-part counts
then yields \((2+o(1))n\) for dense source prime supports below
\(\exp(o(\sqrt{\log M}))\). Exact smooth-part bins also yield \(2n\)
for a separated-band class. Portable certificates combine the older
Hall condition with the new deletion rules. The remaining obstruction
involves a growing prime core and interactions across intermediate
prime scales.
\end{abstract}
\begingroup
\small
\setcounter{tocdepth}{1}
\setlength{\parskip}{1pt}
\makeatletter
\renewcommand{\l@section}[2]{\@dottedtocline{1}{0em}{2em}{#1}{#2}}
\makeatother
\tableofcontents
\endgroup
\newpage

\section{Problem, baseline, and the merged result map}
\label{root:scope}

Let \(A\subseteq\{2,\ldots,M\}\) consist of \(n\) distinct integers,
and let
\[
 I_x=\{x+1,\ldots,x+M\},\qquad x\in\mathbb Z_{\ge0}.
\]
Write \(\tau(A,I_x)\) for the smallest cardinality of a subset
\(B\subseteq I_x\) such that \(\prod_{a\in A}a\mid\prod_{b\in B}b\).
The original problem takes \(M=\max A\). In deletion arguments the
ambient \(M\) is kept fixed, even if the maximum source is deleted.
The global extremal function is
\[
 g(n)=\sup_{|A|=n,\,x\ge0}\tau(A,I_x),\qquad M=\max A.
\]
The targets are \(g(n)\le2n\), or at least
\(g(n)\le(2+o(1))n\). A target position can be used only once.

For every positive integer \(d\),
\begin{equation}
 \left\lfloor\frac Md\right\rfloor
 \le N_I(d):=\#(I_x\cap d\mathbb Z)
 \le\left\lceil\frac Md\right\rceil.
 \label{root:counts}
\end{equation}
The elementary full-supply identity is
\begin{equation}
 \prod A\mid M!\mid\prod I_x,
 \qquad \prod I_x=M!\binom{x+M}{M}.
 \label{root:supply}
\end{equation}
Consequently, whenever a selected product is deficient in a prime
valuation, an unused target containing that prime exists. Adding it
reduces the sum of the positive valuation deficits by at least one.
This repair principle never reuses a selected position.

\subsection{The quadratic range is an external input}

Chen's public work includes the sufficient range \(M\ge8n^3\)
for an exact \(2n\) cover \cite{Chen}. The submitted argument at
\(M\ge32n^{5/2}\) is reconstructed completely below. However, Wang's
written appendix already proves
\begin{equation}
 M\ge16n^2\quad\Longrightarrow\quad\tau(A,I_x)\le2n.
 \label{root:wang}
\end{equation}
We checked its full proof, including the difficult-divisor Hall
argument and the shared position budget \cite[Appendix A]{Wang}.
Thus the quadratic order is not a new conclusion of this synthesis.
The further user-supplied summaries asserted a quadratic threshold
but did not include its complete proof; we use the checked explicit
theorem \eqref{root:wang} instead of treating those summaries as a proof.

At the other end, \(M\le2n\) is immediate by selecting the whole
interval. The unresolved unrestricted exact range is therefore
\begin{equation}
 \boxed{2n<M<16n^2.}
 \label{root:range}
\end{equation}
Wang's work also supplies a uniform linear upper bound and a
fractional-to-integer transfer theorem; we identify explicitly below
where these external inputs are used. They are public working
results, not claimed here as externally refereed results.

\subsection{What is established in this package}

\begin{center}
\small
\begin{tabular}{p{0.43\linewidth}p{0.47\linewidth}}
\toprule
Class or reduction & Conclusion and scope\\
\midrule
All source sets with \(M\ge16n^2\) & Exact \(2n\); external theorem \cite{Wang}\\
Submitted \(M\ge32n^{5/2}\) argument & Reconstructed; its three-large-factor step has a separate range condition\\
Finite cofactor criterion & Exact \(2n\), for every shift; retained from the earlier note\\
All source primes \(\ge M^\alpha\), fixed \(\alpha>1/(2e)\) & Exact \(2n\) for sufficiently large \(M\); retained and reproved\\
All source primes \(\le Y(M)\), \(\log Y=o(\sqrt{\log M})\) & \((2+o(1))n\), uniformly over shifts; proved below\\
One separated large prime per source & Exact \(2n\), allowing a growing dense small-prime support; proved below\\
Arbitrary support, sufficiently small individual raw weights & Uniform \(H_{2+\varepsilon}\), even with one arbitrary exceptional weight\\
At most \(k\) weights exceeding \(c_\varepsilon/(k+1)\) & Explicit hinge defect at most \(2^{k+1}(1+\log_2 M)^{k+1}\)\\
Minimal counterexample in the remaining range & A core of size \(O(\sqrt{n/\log n})\), with at most one outside-prime occurrence per source\\
Unified finite recognizer & Exact all-shift certificates for inputs satisfying a sequence of sufficient rules; otherwise UNKNOWN\\
\bottomrule
\end{tabular}
\end{center}

The earlier full working note is included in the companion bundle
\cite{Earlier}. It also records fixed-support, reciprocal-mass, and
fixed-anchor results. The current note supplies complete proofs for
its principal additional claims. Computational searches are not used
to establish a universal mathematical assertion.

\section{Retained cofactor matching and rough-source covers}
\label{root:hall}

For a composite \(a\in A\), let \(p(a)=P^+(a)\) be its largest
prime factor and put \(d(a)=a/p(a)\). Let \(Q\) be the union of the
prime supports of these cofactors, and let
\[
 r_A(d)=\#\{a\in A:a\text{ composite},\ d(a)=d\}.
\]

\begin{theorem}[Finite all-shift Hall criterion]
\label{root:hall-theorem}
If, for every distinct composite-source cofactor \(d\),
\begin{equation}
 \left\lfloor\frac Md\right\rfloor
 -\sum_{\substack{q\in Q\\dq\le M}}
       \left\lceil\frac M{dq}\right\rceil
 \ge r_A(d),
 \label{root:hall-finite}
\end{equation}
then \(\tau(A,I_x)\le2n\) for every shift \(x\ge0\).
\end{theorem}

\begin{proof}
Give a composite source all target neighbors divisible by its
cofactor. For a subset of the composite sources let \(D\) be its set
of distinct cofactors. Each \(d\in D\) is \(Q\)-smooth and
\(dP^+(d)\le M\). Define
\[
 K_D(b)=\sum_{d\in D}
 \left(\mathbf1_{d\mid b}
       -\sum_{\substack{q\in Q\\q\le M/d}}
            \mathbf1_{dq\mid b}\right).
\]
At most one summand can be positive. Indeed, a positive term forces
\(d\) to be the full smooth part of \(b\) over the primes
\(Q\cap[2,M/d]\). If \(d_1<d_2\) both had this property, the larger
prime set associated with \(d_1\) would have the smaller smooth part,
a contradiction. A positive term has value one, and all other terms
are nonpositive. Outside the union of the neighborhoods every term
vanishes. Hence \(K_D\) is bounded above by the indicator of that
union. Summing and using \eqref{root:counts} and
\eqref{root:hall-finite} gives a union size at least
\(\sum_{d\in D}r_A(d)\), and thus at least the size of the chosen
source subset. Hall's theorem provides distinct cofactor carriers.

Prime sources, whose quotient is one, may be assigned afterward to
unused positions. There is room because \(n\le M-1\). The selected
\(n\) positions supply the product of all cofactors. At most \(n\)
prime units remain missing; \eqref{root:supply} allows the same
number of unused-position repairs. The total is at most \(2n\).
\end{proof}

\begin{corollary}[Rough sources]
\label{root:rough}
Fix \(\alpha>1/(2e)\). If every prime divisor of every source is at
least \(M^\alpha\), then \(\tau(A,I_x)\le2n\) for every shift, once
\(M\) is sufficiently large depending on \(\alpha\).
\end{corollary}

\begin{proof}
Every cofactor prime \(q\) satisfies \(q\le P^+(a)\) and
\(qP^+(a)\mid a\), so \(q\le\sqrt M\). When
\(1/(2e)<\alpha<1/2\), Mertens' prime reciprocal estimate gives
\[
 \sum_{q\in Q}\frac1q
 \le\log\frac1{2\alpha}+o(1)<1.
\]
For \(T=M/d\), each cofactor has \(T\ge P^+(a)\ge M^\alpha\), and
its fibre contains at most \(\pi(T)\) sources. The left side of
\eqref{root:hall-finite} is at least
\[
 T\left(1-\sum_{q\in Q,\,q\le T}\frac1q\right)-\pi(T)-1.
\]
Since \(\sup_{T\ge M^\alpha}(2\pi(T)+1)/T=o(1)\), this is at
least \(\pi(T)\) for sufficiently large \(M\). Apply the theorem.
For \(\alpha>1/2\) all sources are prime; at \(\alpha=1/2\) the
only possible composite boundary is a square and the same finite
criterion applies. The standard prime estimates used here follow
from Mertens' theorem and the elementary prime-counting upper bound
\cite{TaoMertens}.
\end{proof}

\section{Deletion banks and necessary counterexample structure}
\label{sec:comb-banks}

Throughout this section, let $A\subseteq\{2,\ldots,M\}$ be a set of
$n$ distinct integers and let
\[
 I=\{x+1,\ldots,x+M\},\qquad x\in\mathbb Z_{\ge0}.
\]
The ambient length $M$ is retained when sources are deleted; it need
not equal the maximum of a smaller source set. Write
\[
 D_p(A)=\sum_{a\in A}v_p(a),\qquad
 N_I(d)=\#\{b\in I:d\mid b\},
\]
and let $\tau(A,I)$ be the minimum cardinality of a subset whose
product is divisible by $\prod_{a\in A}a$. We use $P^+(a)$ for the
largest prime factor, $\omega(a)$ for the number of distinct prime
factors, and $\Omega(a)$ for their number with multiplicity.

The following two facts will be used repeatedly:
\begin{equation}
 \left\lfloor\frac Md\right\rfloor\le N_I(d)
 \le\left\lceil\frac Md\right\rceil,
 \qquad
 \prod_{a\in A}a\mid M!\mid\prod_{b\in I}b.
 \label{eq:comb-supply}
\end{equation}
The divisibility $\prod A\mid M!$ uses distinctness of the sources;
$M!\mid\prod I$ follows from $\prod I=M!\binom{x+M}{M}$.
In particular, if a chosen subset has deficient $p$-valuation, some
unused interval position contains $p$. Adding that position reduces
the sum of the positive valuation deficits by at least one. We call
this a \emph{prime-unit repair}; different repairs always use
different positions.

\subsection{The local deletion hypothesis}

Let $K$ be a positive integer. For the necessary conditions below,
assume
\begin{equation}
 \tau(A,I)>K+1,
 \qquad
 \text{each }A\setminus\{a\}\text{ has a cover }B_a\subseteq I
 \text{ with }|B_a|\le K-1.
 \label{eq:comb-local}
\end{equation}
All conclusions apply to \emph{every} such cover $B_a$. For a
least-cardinality counterexample to the exact $2n$ bound, one takes
$K=2n-1$. A different choice needed for the asymptotic problem is
justified separately below.

\begin{lemma}[Forced deletion banks and a divisor gap]
\label{lem:comb-forcing}
Under \eqref{eq:comb-local}, for every $a\in A$ and every prime $q\mid a$,
\[
 C_q(a):=\{b\in I:a/q\mid b\}\subseteq B_a.
\]
Consequently every source prime is strictly smaller than $K$,
every source has $\Omega(a)\ge3$, and
\begin{equation}
 a>\frac MK,
 \qquad \frac aq>\frac MK\quad(q\mid a\text{ prime}).
 \label{eq:comb-quotient-size}
\end{equation}
No source has a divisor $d$ with
\begin{equation}
 K<d\le\frac MK.
 \label{eq:comb-divisor-gap}
\end{equation}
\end{lemma}

\begin{proof}
If $C_q(a)$ contained an unused position, add it to $B_a$. All
demands except possibly one $q$-unit are then supplied. One repair
finishes the cover in at most $K+1$ positions, a contradiction.
Since
$|C_q(a)|\ge\lfloor M/(a/q)\rfloor\ge q$, this inclusion gives
$q\le K-1$. It also gives $a/q>M/K$: otherwise the bank would have
at least $K$ members. An unused multiple of the whole source $a$
would finish the cover in one step, proving the analogous bound for
$a$. If $\Omega(a)\le2$, at most two prime-unit repairs would suffice,
which proves the remaining assertion about $\Omega$.

Finally, suppose $a=dv$ with $K<d\le M/K$. There are at least $K$
multiples of $d$, so one lies outside $B_a$. After choosing it, at
most $K$ positions are used. The number of multiples of $v$ is at
least $\lfloor M/v\rfloor\ge d\ge K+1$, so a second fresh position
supplies $v$. Their product supplies $a$, even if $d$ and $v$ are
not coprime, again contradicting \eqref{eq:comb-local}.
\end{proof}

The left endpoint in \eqref{eq:comb-divisor-gap} is strict. In
particular, this gap is empty when $M\le K^2$; it must not be used
as a nontrivial restriction in that range.

\subsection{Two-prime removal and a sufficient source class}

\begin{lemma}[A two-prime demand constraint]
\label{lem:comb-pair}
Assume \eqref{eq:comb-local}. Let $p,q$ be two prime occurrences in
$a\in A$, allowing $p=q$ only when $p^2\mid a$. If $pq\ge K$, then
\begin{equation}
 \begin{split}
 D_p(A)&>v_p\!\left(\prod_{b\in B_a}b\right)
 \ge v_p\!\left(\prod_{b\in C_q(a)}b\right)\\
 &\ge v_p(a/q)\left\lfloor\frac{Mq}{a}\right\rfloor
 \ge q\,v_p(a/q).
 \end{split}
 \label{eq:comb-pair-demand}
\end{equation}
\end{lemma}

\begin{proof}
Suppose instead that $B_a$ already supplies the full $p$-demand of
$A$. There are at least $pq\ge K$ multiples of $a/(pq)$ in $I$,
whereas $|B_a|\le K-1$. Choose one fresh multiple. If $p\ne q$,
every demand is now supplied except possibly one $q$-unit: the
$p$-demand was already complete, and the new position supplies the
other factors of $a$. One repair completes a cover in at most
$K+1$ positions. If $p=q$, the full $p$-demand was already complete
and even the last repair is unnecessary. This contradiction proves
the strict inequality. The others follow from the forced-bank
inclusion and \eqref{eq:comb-supply}.
\end{proof}

Bank inclusion alone does not imply a lower bound on the full
$p$-demand: selected positions could oversupply $p$ while serving
other primes. The fresh multiple of $a/(pq)$ is the additional
argument that makes \eqref{eq:comb-pair-demand} valid.

\begin{theorem}[A sufficient class at every interval scale]
\label{thm:comb-recognition}
Suppose that each $a\in A$ either has $\Omega(a)\le2$, or contains
two prime occurrences $p,q$ for which
\begin{equation}
 pq\ge2n-1,
 \qquad D_p(A)\le q\,v_p(a/q).
 \label{eq:comb-recognition}
\end{equation}
Then $\tau(A,I)\le2n$ for every positive interval of length
$M\ge\max A$. No bound on $\Omega(a)$ is required in the second case.
\end{theorem}

\begin{proof}
Induct on $n$, keeping the ambient length fixed. The case $n=1$
follows by choosing a multiple of its source. The hypothesis is
hereditary: deleting sources decreases both the demands and the
threshold $2n-1$. Obtain a cover $B$ of $A\setminus\{a\}$ using
at most $2n-2$ positions. If $\Omega(a)\le2$, use prime-unit repairs.

Otherwise choose witnesses $p,q$ in \eqref{eq:comb-recognition}.
If $C_q(a)$ has an unused position, add it and repair at most one
$q$-unit. If $C_q(a)\subseteq B$, then
$v_p(\prod B)\ge qv_p(a/q)\ge D_p(A)$. A fresh multiple of
$a/(pq)$ exists, and the construction in Lemma~\ref{lem:comb-pair}
again finishes in at most two additions.
\end{proof}

\subsection{All forced banks must fit simultaneously}

\begin{proposition}[Forced unions and their valuations]
\label{prop:comb-union}
Under \eqref{eq:comb-local}, put
\[
 U_a=\bigcup_{\substack{q\mid a\\q\ \mathrm{prime}}}C_q(a),
 \quad c(a)=1+\sum_{\substack{q\mid a\\q\ \mathrm{prime}}}(q-1),
 \quad t=\left\lfloor\frac Ma\right\rfloor.
\]
Then
\begin{equation}
 |U_a|=\sum_{\substack{q\mid a\\q\ \mathrm{prime}}}N_I(a/q)
       -(\omega(a)-1)N_I(a)\le K-1,
 \qquad t\,c(a)\le K-1.
 \label{eq:comb-union-size}
\end{equation}
For $p\mid a$, write $e_p=v_p(a)$ and define
\begin{equation}
 F_p(a)=t\bigl(e_pc(a)-(p-1)\bigr).
 \label{eq:comb-bank-valuation}
\end{equation}
Then $v_p(\prod U_a)\ge F_p(a)$, and necessarily
\begin{equation}
 \sum_{\substack{p\mid a\\p\ \mathrm{prime}}}
       \min\{e_p,(D_p(A)-F_p(a))_+\}\ge3.
 \label{eq:comb-residual-units}
\end{equation}
If $\Omega(a)=3$, the stronger equalities and inequalities
\begin{equation}
 v_p\!\left(\prod B_a\right)=D_p(A)-e_p,
 \qquad D_p(A)\ge F_p(a)+e_p\quad(p\mid a)
 \label{eq:comb-three-units}
\end{equation}
hold for every deletion cover $B_a$.
\end{proposition}

\begin{proof}
For distinct primes $p,q\mid a$, one has
$\operatorname{lcm}(a/p,a/q)=a$. Thus two different banks intersect
exactly in the multiples of $a$. A multiple of $a$ is counted
$\omega(a)$ times in the sum, and every other union member only
once. This proves the identity in \eqref{eq:comb-union-size}.

In a complete period of length $a$, there is the common residue
$0$, and for each prime $q\mid a$ there are $q-1$ additional
residues $ja/q$, $1\le j<q$. They are disjoint across the primes.
The period therefore contains $c(a)$ union members. Partitioning
the interval into $t$ complete periods and a remainder proves the
uniform cardinality bound.

For a fixed $p\mid a$, the $p-1$ nonzero residues from its own bank
have valuation exactly $e_p-1$. Every other union residue has
valuation at least $e_p$. This proves
\eqref{eq:comb-bank-valuation} on each complete period, and hence
the asserted lower bound on the whole interval.

The deletion cover supplies at least $D_p(A)-e_p$ units at every
prime of $a$, and it contains $U_a$. Its remaining $p$-deficit is
therefore at most $\min\{e_p,(D_p(A)-F_p(a))_+\}$. A sum at most
two would permit two repairs, proving \eqref{eq:comb-residual-units}.
When $\Omega(a)=3$, any surplus over $D_p(A)-e_p$ at even one prime
of $a$ would leave at most two units. Thus all those surpluses are
zero, which gives \eqref{eq:comb-three-units}.
\end{proof}

For example, for a squarefree triple $a=pqr$ the last bound gives
$D_p(A)\ge\lfloor M/a\rfloor(q+r-1)+1$, and its cyclic analogues.
These are necessary conditions, not a proof that such a deletion
cover exists.

\subsection{Saturation by a prime-power bank}

\begin{lemma}[Prime-power saturation]
\label{lem:comb-power-saturation}
Let $A$ consist of $n$ distinct sources in $[2,M]$. If $B\subseteq I$
contains every multiple of $p^e$, where $e\ge1$, and
$\lfloor M/p^e\rfloor\ge n$, then $v_p(\prod B)\ge D_p(A)$.
\end{lemma}

\begin{proof}
Counting source valuation levels gives
\[
 D_p(A)\le\sum_{j\ge1}\min\left\{n,
                         \left\lfloor\frac M{p^j}\right\rfloor\right\}.
\]
Every multiple of $p^e$ contributes its first $e$ units, and all
interval multiples of higher powers belong to the same bank. Hence
\[
 \begin{split}
 v_p\!\left(\prod B\right)
 &\ge eN_I(p^e)+\sum_{j>e}N_I(p^j)\\
 &\ge en+\sum_{j>e}\left\lfloor\frac M{p^j}\right\rfloor
 \ge D_p(A).
 \end{split}
\]
\end{proof}

\begin{corollary}
\label{cor:comb-power-exclusion}
Under \eqref{eq:comb-local}, no source has the form
$a=p^e q$, with $p,q$ prime, $e\ge1$, and $p^e\le M/n$.
The primes $p,q$ may coincide. In particular, such a source is
excluded whenever $q\ge n$.
\end{corollary}

\begin{proof}
Deleting the $q$ occurrence forces all multiples of $p^e$ into
$B_a$. Lemma~\ref{lem:comb-power-saturation} supplies the full
$p$-demand. Only one $q$-unit can remain; if $q=p$, nothing remains.
This contradicts \eqref{eq:comb-local}. Finally,
$q\ge n$ implies $p^e=a/q\le M/n$.
\end{proof}

The prime-power form is essential to this proof. It does not replace
the threshold $M/K$ by $M/n$ for an arbitrary composite cofactor.

\subsection{A small core and one remaining prime per source}

\begin{theorem}[Structure of a deletion obstruction]
\label{thm:comb-core}
Assume \eqref{eq:comb-local}. Among the source primes, define
\[
 E=\{p:p\le\sqrt K\}\ \cup\ \{p:D_p(A)>p\},
 \qquad R=\{\text{source primes}\}\setminus E.
\]
Every $r\in R$ is the unique largest prime of any source containing
it, and occurs there to exponent one. Thus such a source is $a=rd$,
where $d$ is $E$-smooth and $P^+(d)<r$, with $P^+(1)=1$.
Moreover
\begin{equation}
 |E|\le\lfloor\sqrt K\rfloor+\sqrt{2n\log_2M},
 \qquad R\subset(\sqrt K,K).
 \label{eq:comb-core-elementary}
\end{equation}
If $n\le K\le C_0n$ and $M<16n^2$, with fixed $C_0$, then as
$n\to\infty$,
\begin{equation}
 |E|=O_{C_0}\!\left(\sqrt{\frac n{\log n}}\right),\qquad
 \sum_{r\in R}\frac1r\le\log2+o(1),\qquad
 \sum_{\substack{p\in E\\p>\sqrt K}}\frac1p=o(1).
 \label{eq:comb-core-sharp}
\end{equation}
\end{theorem}

\begin{proof}
If a source containing $r\in R$ had another prime occurrence
$q\ge r$, then $rq>K$ and
$D_r(A)\le r\le qv_r(a/q)$, contradicting
Lemma~\ref{lem:comb-pair}. This proves both the unique-largest-prime
assertion and the fact that the remaining cofactor is $E$-smooth.
Lemma~\ref{lem:comb-forcing} gives the upper endpoint $r<K$.

Let $H=\{p:D_p(A)>p\}$ and $h=|H|$. Distinctness of its primes gives
\[
 \frac{h(h+1)}2\le\sum_{p\in H}p
 <\sum_pD_p(A)=\sum_{a\in A}\Omega(a)\le n\log_2M,
\]
which proves the elementary bound. For the sharper estimate use
the logarithmic budget
\begin{equation}
 \sum_{p\in H}p\log p
 <\sum_pD_p(A)\log p
 =\sum_{a\in A}\log a\le n\log M.
 \label{eq:comb-log-budget}
\end{equation}
Chebyshev's bound $\pi(t)\ll t/\log t$ implies that the $j$th prime
is at least a constant times $j\log(j+1)$. For completeness, the
usual bound follows by applying
$\prod_{m<p\le2m}p\mid\binom{2m}{m}<4^m$, summing over dyadic
intervals to bound $\sum_{p\le t}\log p$, and separating primes
below and above $\sqrt t$. It follows that the sum of $p\log p$
over any $h$ primes is at least
$c_1h^2\log^2(h+1)$ for an absolute $c_1>0$.
Together with \eqref{eq:comb-log-budget} and $\log M=O(\log n)$,
this gives $h=O(\sqrt{n/\log n})$: if $h<n^{1/4}$ it is immediate,
and otherwise $\log(h+1)\ge\frac14\log n$.
The remaining part of $E$ has at most
$\pi(\sqrt K)=O_{C_0}(\sqrt n/\log n)$ primes.

The prime reciprocal asymptotic
$\sum_{p\le t}1/p=\log\log t+B+o(1)$, for example from
Mertens' estimates \cite{TaoMertens}, gives
\[
 \sum_{\sqrt K<p<K}\frac1p=\log2+o(1).
\]
To control the large-prime part of $E$, put $T=\sqrt K\log n$.
The sum over all primes in $(\sqrt K,T]$ tends to zero because
$\log T/\log\sqrt K\to1$. For the rest, use
\eqref{eq:comb-log-budget}:
\[
 \sum_{\substack{p\in H\\p>T}}\frac1p
 \le\frac{n\log M}{T^2\log T}=O_{C_0}(\log^{-2}n).
\]
This proves all of \eqref{eq:comb-core-sharp}.
\end{proof}

The set $E$ is defined by source demands, not by the magnitudes of
dual weights. Its sublinear size does not by itself supply a
fractional cover or a rounding argument for its interaction with $R$.

\subsection{Why the reduction also applies to an asymptotic obstruction}

Let $g(n)$ denote the supremum of $\tau(A,I)$ over $n$ distinct
sources, with interval length $M=\max A$. A smallest exact-$2n$
counterexample does not automatically describe failure of the
asymptotic conjecture. The following argument provides the required
quantifiers.

\begin{proposition}[Linear budgets with an additive constant]
\label{prop:comb-asym-minimality}
Fix $2<c<10$. Suppose that for every $C\ge0$ there is an instance
with $\tau(A,I)>cn+C$. For each $C$, choose the smallest failing
cardinality $n_C$ and one such instance. As $C\to\infty$, these
instances satisfy the conclusions of Theorem~\ref{thm:comb-core}
with
\[
 K=\lfloor c(n_C-1)+C\rfloor+1,
 \qquad 2n_C-1\le K<10n_C,
 \qquad M<16n_C^2,
\]
and $n_C\to\infty$.
\end{proposition}

\begin{proof}
Minimality supplies a cover of every deletion with at most
$\lfloor c(n_C-1)+C\rfloor=K-1$ positions. If the deletion has
smaller maximum, apply its bound inside a contained interval of
that smaller length. Since $c\ge2$,
\[
 K+1\le\lfloor cn_C+C\rfloor<\tau(A,I).
\]
Thus \eqref{eq:comb-local} holds.

The known uniform bound $g(n)\le10n$ \cite{Wang} implies
$C<(10-c)n_C$, hence $n_C\to\infty$ and $K<10n_C$; the lower
bound on $K$ follows from $c\ge2$ and $C\ge0$.
The known long-interval theorem
$M\ge16n_C^2\Rightarrow\tau(A,I)\le2n_C$ \cite{Wang}
forces $M<16n_C^2$. Theorem~\ref{thm:comb-core} now applies.
\end{proof}

If the unrestricted asymptotic coefficient-two assertion fails,
one can choose a fixed $c>2$ below the offending limiting upper
ratio, and then no finite additive constant $C$ works. Thus the
proposition is a genuine necessary reduction for a possible
asymptotic obstruction family. It does not prove that the growing
core can be absorbed, nor identify it with the exceptional support
of an optimal dual vector.

\section{A complete reconstruction of the \texorpdfstring{$32n^{5/2}$}{32 n to the five-halves} argument}
\label{sec:comb-dot}

This section reconstructs the supplied argument in full, including
the demand step and the parameter restriction on which its
three-occurrence structure depends. Its range is already contained
in the stronger quadratic result of Wang \cite{Wang}; it is not
presented as a new best long-interval bound.

\begin{theorem}
\label{thm:comb-dot}
For every set $A$ of $n$ distinct integers in $[2,M]$ and every
positive interval $I$ of length $M$,
\[
 M\ge32n^{5/2}\quad\Longrightarrow\quad\tau(A,I)\le2n.
\]
\end{theorem}

\begin{proof}
Induct on $n$ at fixed ambient length $M$. The case $n=1$ is
immediate. The induction hypothesis is available for a deletion
because $M\ge32n^{5/2}\ge32(n-1)^{5/2}$. Suppose a counterexample
exists with $n\ge2$, and set
\[
 K=2n-1,\qquad c=\frac M{K^3},\qquad L=\frac MK=cK^2.
\]
The local deletion hypothesis holds. In particular, every source
prime is below $K$, every prime-deletion quotient exceeds $L$,
and no source divisor lies in $(K,L]$.

\paragraph{The parameter range.}
If $c\ge1$, every source exceeds $L\ge K^2$. Its least divisor
$d>K$ has all proper divisors at most $K$. Removing any prime
factor $p<K$ gives $d/p\le K$, so $d<K^2\le L$, a contradiction.
We may therefore assume $0<c<1$. The numerical hypothesis gives
\begin{equation}
 c\ge\frac{32n^{5/2}}{(2n-1)^3}>\frac4{\sqrt n},
 \qquad c^2K\ge1.
 \label{eq:comb-dot-range}
\end{equation}
All structural steps that follow use the explicit restrictions
$0<c<1$ and $c^2K\ge1$. They imply
$cK>1$ and $c^{3/2}K>1$.

\paragraph{Finding a large prime occurrence.}
Choose the least divisor $d>K$ of a source. If $\Omega(d)\ge3$,
its smallest prime factor is at most $d^{1/3}$; removing it leaves
a proper divisor at least $d^{2/3}$. Thus
\[
 d\le K^{3/2}\le cK^2=L,
\]
contradicting the divisor gap. The divisor is not prime because
all source primes are below $K$. Hence $d$ is a product of two
prime occurrences, and its larger prime $q$ satisfies
\begin{equation}
 q>\sqrt c\,K.
 \label{eq:comb-dot-large-prime}
\end{equation}

\paragraph{The small part.}
Put $U=1/\sqrt c$. There is no source divisor $t$ with
$U<t\le cK$. Indeed, the prime in
\eqref{eq:comb-dot-large-prime} is larger than $t$, since
$\sqrt cK>cK$. Therefore $qt$ divides the source and
\[
 K<qt<cK^2=L,
\]
again impossible. Every source prime is consequently either at
most $U$ or larger than $cK$.

Let $u$ be the full part of a source supported on primes at most
$U$. Then $u\le U$. Otherwise its least divisor $t>U$ is at most
$U^2=1/c$: remove a prime factor at most $U$ and use minimality.
But $1/c\le cK$ by \eqref{eq:comb-dot-range}, contradicting the
preceding forbidden interval.

\paragraph{Exactly three large occurrences.}
Call a prime occurrence larger than $cK$ large. Any two large
occurrences have product greater than $c^2K^2\ge K$. The divisor
gap therefore forces their product above $L$. Four such
occurrences would have product greater than
$L^2=M(cK)>M$, so a source has at most three.

It cannot have at most two. Delete a prime occurrence so that at
most one large occurrence remains; if none is large, delete any
prime. The quotient is then at most
\[
 uK\le UK=\frac K{\sqrt c}<cK^2=L,
\]
where $c^{3/2}K>1$ was used. This contradicts
\eqref{eq:comb-quotient-size}. Thus every source has the canonical
form
\begin{equation}
 a=u p q r,\qquad u\le U,
 \label{eq:comb-dot-triples}
\end{equation}
where $p,q,r$ are three large prime occurrences, with repetition
allowed. The small part $u$ need not have a bounded number of
prime factors.

\paragraph{Full demands at the large primes.}
Fix any large occurrence $p$ in a source
\eqref{eq:comb-dot-triples}. The other two have product above
$L=cK^2$, so one of them, say $q$, is larger than $\sqrt cK$.
Also $pq>L>K$. Lemma~\ref{lem:comb-pair} gives
\begin{equation}
 D_p(A)>qv_p(a/q)\ge q>\sqrt cK.
 \label{eq:comb-dot-demand}
\end{equation}
This remains valid when occurrences have equal prime values.
It is the two-prime removal argument, rather than bank inclusion
alone, that justifies \eqref{eq:comb-dot-demand}.

Let $h$ be the number of distinct large primes across the sources.
There are exactly $3n$ large-prime occurrences in total, so
\[
 h<\frac{3n}{\sqrt cK}\le\frac2{\sqrt c}=2U,
\]
where $3n\le2(2n-1)$ for $n\ge2$.

\paragraph{Counting the distinct sources.}
Unique factorization makes \eqref{eq:comb-dot-triples} an injection
into the pairs consisting of an integer $u\le U$ and a multiset
of three primes from the $h$ available large primes. Therefore
\[
 n\le\lfloor U\rfloor\binom{h+2}{3}
   <\frac{U(4U)^3}{6}
   =\frac{64}{6c^2}
   <\frac23n.
\]
Here $U>1$, $h<2U$, and $c>4/\sqrt n$ were used. The contradiction
completes the induction.
\end{proof}

The same three-occurrence argument cannot simply be continued to
$M=Cn^2$ with fixed $C$. Then $c=M/(2n-1)^3$ is of order $1/n$,
and the essential condition $c^2K\ge1$ fails for large $n$.
A quadratic proof therefore requires a different argument; the
independently checked theorem $M\ge16n^2\Rightarrow\tau\le2n$
\cite{Wang} supplies that range. None of the necessary conditions
above resolves the remaining unrestricted exact or asymptotic
problem.

\section{Diffuse weights and dense prime supports}
\label{an:section}

This section proves three uniform estimates for the raw hinge formulation of
the interval divisor-cover problem. The first allows an arbitrary prime
support when all weights are sufficiently small, together with an extension
allowing one arbitrary exceptional prime weight. The second allows a
controlled number of exceptional weights, at the cost of an explicit
additive defect. The third treats every weight vector on primes below
\(\exp(o(\sqrt{\log M}))\), and gives an asymptotic coefficient-two
integer cover for source sets supported on those primes. These results do
not establish the unrestricted conjecture.

All logarithms are natural unless a base is indicated, and
\((t)_+=\max(t,0)\). Prime sets in this section are finite. For nonnegative
weights \(w=(w_p)_{p\in P}\), write
\[
 X_w(m)=\sum_{p\in P}w_pv_p(m).
\]
If \(I\) consists of \(M\) consecutive positive integers, define
\[
 \Delta_c(M,I,w)
 =\sum_{a=1}^{M}(X_w(a)-c)_+
   -\sum_{b\in I}(X_w(b)-1)_+.
\]
The assertion \(H_c\) is \(\Delta_c(M,I,w)\le0\). Whenever coprimality
counts at real lengths are used, an interval of length \(T\) means a
half-open interval \((x,x+T]\), with its integer points understood and
\(x\ge0\). For a finite prime set \(Q\), put
\[
 Q^{\Pi}=\prod_{q\in Q}q,
 \qquad \vartheta(Q)=\prod_{q\in Q}(1-1/q).
\]
Both empty products equal one. An integer is \(Q\)-smooth if every prime
divisor belongs to \(Q\); the integer one is included.

\subsection{A uniform consequence of the fundamental lemma}

\begin{lemma}[Uniform sifting at a large parameter]
\label{an:sieve}
There is an absolute constant \(C>0\) such that the decreasing function
\begin{equation}
 e(t)=C\bigl((t/2)^{-t/4}+2^{1-t/2}\bigr),\qquad t\ge4,
 \label{an:base-error}
\end{equation}
has the following property. If \(z\ge2\), \(Q\) is any subset of the
primes at most \(z\), and \(J\) is an interval of real length
\(T\ge z^t\), then
\[
 (1-e(t))T\vartheta(Q)
 \le \#\{m\in J:(m,Q^{\Pi})=1\}
 \le (1+e(t))T\vartheta(Q).
\]
In particular, the estimate is uniform in the interval shift and in the
choice of the subset \(Q\), and \(e(t)\to0\) as \(t\to\infty\).
\end{lemma}

\begin{proof}
The fundamental lemma of the combinatorial sieve, in dimension one, gives
\[
 \#\{m\in J:(m,Q^{\Pi})=1\}
 =T\vartheta(Q)\{1+O(u^{-u/2})\}+O(z^u),\qquad u\ge1.
\]
The constants are absolute here. Indeed, the remainder in each
divisibility count is at most one, while the dimension-one product
condition holds with common constants for every subset of the primes.
This is the fundamental-lemma statement and interval application in
\cite{Lagarias}; see also the interval version for subsets of the small
primes in \cite[Theorem~1.4.2]{GranvilleSound}.

Take \(u=t/2\). The elementary product estimate
\[
 \vartheta(Q)
 \ge \prod_{m=2}^{\lfloor z\rfloor}(1-1/m)
 =\frac{1}{\lfloor z\rfloor}\ge\frac1z
\]
shows that the relative contribution of the additive error is at most
\[
 O\!\left(\frac{z^{t/2}}{T\vartheta(Q)}\right)
 =O(z^{1-t/2})=O(2^{1-t/2}).
\]
Increasing a single absolute constant gives the stated function. Both
terms in that function decrease for \(t\ge4\).
\end{proof}

\subsection{A level comparison for diffuse raw weights}

The next proposition works directly with the unbounded valuations
\(v_p\). The final prime-power block of a carrier is allowed to be
partial; this will be essential for controlling large exponents.

\begin{proposition}[Diffuse level comparison]
\label{an:diffuse-level}
Let \(\ell>\eta>0\), let \(0\le w_p\le\eta\), and put
\[
 \mu=\sum_{p\in P}\frac{w_p}{p-1},\qquad
 r=\lfloor\ell/\eta\rfloor\ge8,\qquad
 j_0=\lfloor r/2\rfloor,\qquad \kappa=r-j_0.
\]
Suppose
\begin{equation}
 (1+e(\kappa))\mu+\eta2^{1-j_0}
 \le(\ell-\eta)(1-e(r)).
 \label{an:diffuse-condition}
\end{equation}
Then, for every real \(T>0\), every interval \(J\) of length \(T\), and
every \(s\ge0\),
\begin{equation}
 \#\{a\le T:X_w(a)>s+\ell\}
 \le \#\{b\in J:X_w(b)>s\}.
 \label{an:diffuse-level-ineq}
\end{equation}
Consequently \(H_{1+\ell}\) holds at every integer length and shift.
\end{proposition}

\begin{proof}
Order the prime occurrences from \(P\) in a source integer \(a\) by
increasing prime, retaining multiplicity and omitting primes outside
\(P\). If \(X_w(a)>s+\ell\), let \(d\) be the product of the first
occurrences up to the first crossing of \(s\), and let \(p\) be its
largest prime divisor. Then \(d>1\) and
\begin{equation}
 X_w(d/p)\le s<X_w(d)\le s+\eta.
 \label{an:partial-crossing}
\end{equation}
All \(P\)-primes below \(p\) occur in \(d\) to their full exponents
in \(a\). Further copies of \(p\) may remain in \(a/d\). In particular,
\[
 X_w(a/d)>\ell-\eta.
\]
Every prime occurrence contributes at most \(\eta\), so \(a/d\) has at
least \(r\) occurrences from \(P\), counting multiplicity. All of them
are at least \(p\). Thus \(a/d\ge p^r\), and \(T/d\ge p^r\).

Let \(\mathcal F_s\) consist of all \(P\)-smooth \(d>1\) satisfying
\eqref{an:partial-crossing} and \(dp^r\le T\), where \(p\) is the
largest prime divisor of \(d\). For such a carrier, put
\[
 Q_p=\{q\in P:q<p\},\qquad V_p=\vartheta(Q_p).
\]
Define \(K_d(b)=0\) when \(d\nmid b\), and otherwise put
\[
 K_d(b)=\mathbf1_{(b/d,Q_p^{\Pi})=1}.
\]
For each positive integer \(b\),
\begin{equation}
 \sum_{d\in\mathcal F_s}K_d(b)
 \le \mathbf1_{X_w(b)>s}.
 \label{an:partial-certificate}
\end{equation}
Indeed, positivity forces the factors of \(b\) at all \(P\)-primes
below \(p\) to be exactly those in \(d\). At \(p\), condition
\eqref{an:partial-crossing} specifies the unique first occurrence
crossing \(s\). Hence at most one term is positive, and it forces
\(X_w(b)>s\).

Write \(U=T/d\). The multiples of \(d\) in \(J\), divided by \(d\),
form the integer points of an interval of length \(U\ge p^r\).
Lemma~\ref{an:sieve} and \eqref{an:partial-certificate} imply
\begin{equation}
 \#\{b\in J:X_w(b)>s\}
 \ge \sum_{d\in\mathcal F_s}(1-e(r))UV_p.
 \label{an:diffuse-target}
\end{equation}

For a fixed carrier, enlarge its possible source cofactors to
\[
 \mathcal A_d=\{t\le U:(t,Q_p^{\Pi})=1, X_w(t)>\ell-\eta\}.
\]
Every member still has at least \(r\) occurrences from \(P\), all at
least \(p\). Therefore
\begin{equation}
 |\mathcal A_d|
 \le\frac1{\ell-\eta}\sum_{t\in\mathcal A_d}X_w(t).
 \label{an:cofactor-markov}
\end{equation}
Expand the sum into prime-power divisibility indicators. If \(q\ge p\),
\(j\le j_0\), and some \(t\in\mathcal A_d\) is divisible by \(q^j\),
then \(t/q^j\) has at least \(r-j\) such occurrences. It follows that
\[
 U/q^j\ge p^{r-j}\ge p^{\kappa}.
\]
Terms for which this length inequality fails are zero before enlarging
the cofactor set. For every other term, dividing by \(q^j\) preserves
coprimality to \(Q_p\). Lemma~\ref{an:sieve} therefore bounds the
contribution of all \(j\le j_0\) by
\begin{equation}
 (1+e(\kappa))UV_p
   \sum_{q\in P,\ q\ge p}w_q\sum_{j\ge1}q^{-j}
 \le(1+e(\kappa))UV_p\mu.
 \label{an:low-powers}
\end{equation}

For \(j>j_0\), discard coprimality and the heavy-cofactor restriction.
Their contribution is at most
\[
 U\sum_{q\in P,\ q\ge p}\frac{w_q}{q^{j_0}(q-1)}
 \le U\eta\sum_{m=p}^{\infty}\frac1{m^{j_0}(m-1)}
 \le\frac{U\eta p^{1-j_0}}{p-1}.
\]
The last estimate uses
\(\sum_{m=p}^{\infty}1/[m(m-1)]=1/(p-1)\). Also
\[
 V_p\ge\prod_{m=2}^{p-1}(1-1/m)=\frac1{p-1},
\]
including \(p=2\). Thus the high-power contribution is at most
\begin{equation}
 UV_p\eta2^{1-j_0}.
 \label{an:high-powers}
\end{equation}
Each heavy source is assigned to one carrier, so its total count is at
most \(\sum_d|\mathcal A_d|\). Equations~\eqref{an:cofactor-markov}--\eqref{an:high-powers}, followed
by \eqref{an:diffuse-condition}, show that
\(|\mathcal A_d|\le(1-e(r))UV_p\). Summing over carriers and using
\eqref{an:diffuse-target} proves \eqref{an:diffuse-level-ineq}.
All counts used real interval lengths, so no extra rounding assumption
is needed. Finally, integrating \eqref{an:diffuse-level-ineq} over
\(s\ge1\) proves \(H_{1+\ell}\).
\end{proof}

\subsection{Large weighted means and the diffuse theorem}

\begin{lemma}[Elementary prefix moments]
\label{an:moments}
Suppose \(M\ge2\), \(P\subseteq[2,M]\), and \(0\le w_p\le\eta\).
Let \(\mathbb E_0\) denote averaging over \(1,\ldots,M\), and put
\(\mu=\sum_{p\in P}w_p/(p-1)\). Then
\begin{equation}
 \mu-3\eta\le\mathbb E_0X_w\le\mu,
 \qquad
 \mathbb E_0(X_w-\mu)^2\le8\eta\mu.
 \label{an:moment-bounds}
\end{equation}
If \(0<\delta\le1\), \(\eta\le\delta^2/32\), and
\(\mu\ge1+\delta\), then
\begin{equation}
 \sum_{a=1}^{M}\min(X_w(a),2)\ge M.
 \label{an:capped-mass}
\end{equation}
\end{lemma}

\begin{proof}
For \(e_p=\lfloor\log M/\log p\rfloor\), the rounding loss from
prime powers \(p^j\le M\) is at most \(\eta\): distinct prime powers
are distinct integers, and there are at most \(M-1\) of them. The
omitted tail for each \(p\) is
\[
 \frac{w_p}{p^{e_p}(p-1)}\le\frac{2\eta}{M},
\]
because \(p^{e_p+1}>M\). There are at most \(M\) supporting primes,
so the total tail is at most \(2\eta\). This proves the mean estimate.

Counting multiples gives, for distinct primes \(p,q\),
\[
 \mathbb E_0[v_pv_q]\le\frac1{(p-1)(q-1)},
 \qquad
 \mathbb E_0[v_p^2]\le\sum_{j\ge1}(2j-1)p^{-j}
 =\frac{p+1}{(p-1)^2}.
\]
Consequently
\[
 \mathbb E_0X_w^2
 \le\mu^2+\sum_p\frac{w_p^2p}{(p-1)^2}
 \le\mu^2+2\eta\mu.
\]
Combining this with the lower mean estimate proves the second assertion
in \eqref{an:moment-bounds}.

If \(1+\delta\le\mu\le4\), the map \(x\mapsto\min(x,2)\) is
one-Lipschitz, whence
\[
 \mathbb E_0\min(X_w,2)
 \ge\min(\mu,2)-\sqrt{8\eta\mu}
 \ge1+\delta-\sqrt{32\eta}\ge1.
\]
If \(\mu\ge4\), Chebyshev's inequality gives
\[
 \Pr_0(X_w<2)\le\frac{8\eta\mu}{(\mu-2)^2}\le8\eta,
\]
and therefore
\(\mathbb E_0\min(X_w,2)\ge2(1-8\eta)\ge1\).
This proves \eqref{an:capped-mass}.
\end{proof}

For any nonnegative raw weights, prime-power divisibility counts imply
\(\sum_{b\in I}X_w(b)\ge\sum_{a\le M}X_w(a)\). Thus the elementary
identity \((x-c)_+=x-\min(x,c)\) gives
\begin{equation}
 \Delta_2(M,I,w)
 \le M-\sum_{a=1}^{M}\min(X_w(a),2).
 \label{an:automatic-hinge}
\end{equation}
In particular, \eqref{an:capped-mass} implies \(H_2\). Adding further
nonnegative prime coordinates preserves this sufficient condition.

\begin{theorem}[Uniform diffuse-weight hinge inequality]
\label{an:diffuse-theorem}
For every \(\varepsilon>0\) there exists \(\eta_\varepsilon>0\) such
that \(H_{2+\varepsilon}\) holds at every positive integer length and
every shift for every finite prime support with
\(0\le w_p\le\eta_\varepsilon\).
\end{theorem}

\begin{proof}
It suffices to treat \(0<\varepsilon\le1\), since increasing the prefix
threshold decreases its hinge. Put \(\ell=1+\varepsilon\) and
\(\delta=\varepsilon/2\). Choose \(\eta>0\) sufficiently small that
\(\eta\le\varepsilon^2/128\), \(r=\lfloor\ell/\eta\rfloor\ge8\),
and \eqref{an:diffuse-condition} holds with \(\mu\) replaced by
\(1+\varepsilon/2\). Such a choice exists: as \(\eta\downarrow0\),
the two sides tend respectively to \(1+\varepsilon/2\) and
\(1+\varepsilon\).

The case \(M=1\) is immediate. For \(M\ge2\), remove primes greater
than \(M\) before computing \(\mu\). They contribute nothing to the
prefix, and restoring them only increases the interval hinge. If
\(\mu\le1+\varepsilon/2\), Proposition~\ref{an:diffuse-level} applies.
If \(\mu\ge1+\varepsilon/2\), Lemma~\ref{an:moments} and
\eqref{an:automatic-hinge} give the stronger assertion \(H_2\).
\end{proof}

\begin{corollary}[One arbitrary exceptional prime]
\label{an:one-heavy}
The conclusion of Theorem~\ref{an:diffuse-theorem} remains valid if one
prime \(q\) has an arbitrary nonnegative weight, while every other
weight is at most \(\eta_\varepsilon\). The exceptional weight need
not be at most one in this auxiliary assertion.
\end{corollary}

\begin{proof}
Write \(V\) for the raw load from primes other than \(q\), and first
discard primes greater than \(M\) as above. Let \(\mu_L\) be the mean
parameter of these light primes. If \(\mu_L\ge1+\varepsilon/2\),
Lemma~\ref{an:moments} gives
\(\sum_{a\le M}\min(V(a),2)\ge M\). Adding the \(q\)-coordinate
preserves this inequality, so \eqref{an:automatic-hinge} proves \(H_2\).

Suppose \(\mu_L\le1+\varepsilon/2\), and put
\(\ell=1+\varepsilon\), \(c=1+\ell\). The proof of
Theorem~\ref{an:diffuse-theorem} supplies the level comparison
\eqref{an:diffuse-level-ineq} for \(V\), at every length and shift and
every nonnegative level. It remains available after restricting to any
smaller length; primes above that length may again be removed and then
restored. Also \(V(q^jt)=V(t)\).

First adjoin a nondecreasing profile \(u(j)\in[0,1]\), with \(u(0)=0\).
For \(h=c\) or \(h=1\), the telescoping hinge identity is
\begin{align}
 (V(n)+u(v_q(n))-h)_+
 &= (V(n)-h)_+ \notag\\
 &\quad+\sum_{j\ge1}\mathbf1_{q^j\mid n}
   \int_{u(j-1)}^{u(j)}\mathbf1_{V(n)>h-v}\,dv.
 \label{an:one-heavy-increments}
\end{align}
The base hinges satisfy \(H_c\). In each increment, the target threshold
is \(s=1-v\ge0\) and the source threshold is \(c-v=s+\ell\).
Multiples of \(q^j\) in the prefix rescale to
\(1,\ldots,\lfloor M/q^j\rfloor\); those in \(I\) rescale to a
consecutive interval containing at least that many integers. Apply the
level comparison on a subinterval of this length. If the length is zero,
the prefix term vanishes. Extra target positions are nonnegative.
Thus each prefix increment is bounded by its target increment, proving
\(H_c\) for \(V+u(v_q)\).

For the raw exceptional weight, take
\[
 u(j)=\min(w_qj,1),\qquad R(j)=(w_qj-1)_+.
\]
Pointwise,
\[
 (V+u+R-c)_+\le R+(V+u-c)_+,
 \qquad
 (V+u+R-1)_+=R+(V+u-1)_+.
\]
The identity holds because \(R>0\) forces \(u=1\). The nondecreasing
\(q\)-adic function \(R\) has prefix sum at most its sum over \(I\),
by expanding its nonnegative successive increments and counting
multiples of \(q^j\). This proves the raw assertion.
\end{proof}

\subsection{A controlled number of exceptional weights}

\begin{lemma}[Sifting with arbitrary exceptional primes]
\label{an:exceptional-sieve}
Let \(E\) contain at most \(k\) primes, and let \(Q\) be a set of
primes at most \(z\ge2\), disjoint from \(E\). Put
\begin{equation}
 e_k(t)=3^ke(t/2)+3\cdot4^k2^{1-t/2},\qquad t\ge8.
 \label{an:exceptional-error}
\end{equation}
If an interval \(J\) has real length \(T\ge z^t\), then
\[
 \left|\#\{m\in J:(m,Q^{\Pi}E^{\Pi})=1\}
       -T\vartheta(Q)\vartheta(E)\right|
 \le e_k(t)T\vartheta(Q)\vartheta(E).
\]
The estimate is uniform in the exceptional primes, including primes
larger than \(T\). For each fixed \(k\), \(e_k(t)\) decreases to zero.
\end{lemma}

\begin{proof}
Split \(E=E_{\le z}\cup E_{>z}\), put
\(Q_0=Q\cup E_{\le z}\), \(V_0=\vartheta(Q_0)\), and
\(F=\prod_{p\in E_{>z}}p\). Inclusion--exclusion, with \(\mu(f)\)
denoting the M\"obius function, gives
\[
 \#\{m\in J:(m,Q^{\Pi}E^{\Pi})=1\}
 =\sum_{f\mid F}\mu(f)\#\{m\in J/f:(m,Q_0^{\Pi})=1\}.
\]
For \(f\le T/z^{t/2}\), Lemma~\ref{an:sieve} applies at length
\(T/f\), with relative error \(e(t/2)\). The sum of these errors is
at most \(e(t/2)TV_0\sum_{f\mid F}1/f\). Relative to
\(TV_0\vartheta(E_{>z})\), it is bounded by
\[
 e(t/2)\prod_{p\in E_{>z}}\frac{1+1/p}{1-1/p}
 \le3^ke(t/2).
\]
For each remaining \(f\), its exact count is at most
\(T/f+1\le z^{t/2}+1\), while the magnitude of its main term is at
most \(z^{t/2}\). There are at most \(2^k\) such divisors, so their
combined error is at most \(3\cdot2^kz^{t/2}\). Since
\(V_0\ge1/z\) and \(\vartheta(E_{>z})\ge2^{-k}\), the relative
error is at most
\[
 3\cdot4^kz^{1-t/2}\le3\cdot4^k2^{1-t/2}.
\]
Combining the two bounds proves the lemma.
\end{proof}

\begin{proposition}[Conditional diffuse level comparison]
\label{an:conditional-level}
Let \(E\) contain at most \(k\) primes, and let
\(V(m)=\sum_{p\in L}w_pv_p(m)\), where \(L\cap E=\varnothing\) and
\(0\le w_p\le\eta\). Put
\[
 \mu_L=\sum_{p\in L}\frac{w_p}{p-1},\quad
 r=\lfloor\ell/\eta\rfloor\ge16,\quad
 j_0=\lfloor r/2\rfloor,\quad \kappa=r-j_0,
\]
where \(\ell>\eta\). If
\begin{equation}
 (1+e_k(\kappa))\mu_L+\eta2^{k+1-j_0}
 \le(\ell-\eta)(1-e_k(r)),
 \label{an:conditional-condition}
\end{equation}
then for every real \(T>0\), interval \(J\) of length \(T\), and
\(s\ge0\),
\begin{equation}
 \#\{a\le T:(a,E^{\Pi})=1, V(a)>s+\ell\}
 \le\#\{b\in J:(b,E^{\Pi})=1, V(b)>s\}.
 \label{an:conditional-level-ineq}
\end{equation}
\end{proposition}

\begin{proof}
Use the first crossing in the ordered light-prime occurrences, allowing
a partial last block, as in Proposition~\ref{an:diffuse-level}. Each
carrier \(d>1\), with largest prime \(p\), satisfies
\[
 V(d/p)\le s<V(d)\le s+\eta,\qquad T/d\ge p^r.
\]
Let \(Q_p=\{q\in L:q<p\}\) and \(V_p=\vartheta(Q_p)\). For
\(d\mid b\), use the certificate
\[
 K_d(b)=\mathbf1_{(b/d,E^{\Pi}Q_p^{\Pi})=1},
\]
and put \(K_d(b)=0\) otherwise. A positive term is the unique first
crossing of \(s\) among the light-prime occurrences of \(b\). Since
\(d\) is coprime to \(E^{\Pi}\), positivity also forces
\((b,E^{\Pi})=1\). Thus the sum of the certificates is bounded by the
target indicator in \eqref{an:conditional-level-ineq}.

Write \(U=T/d\). Lemma~\ref{an:exceptional-sieve} gives a target
contribution at least
\((1-e_k(r))UV_p\vartheta(E)\) for each carrier. Enlarge its source
cofactors to those \(t\le U\) coprime to \(E^{\Pi}Q_p^{\Pi}\)
with \(V(t)>\ell-\eta\), and apply the Markov bound
\eqref{an:cofactor-markov}. Such cofactors have at least \(r\) light
prime occurrences, all at least \(p\).

For \(j\le j_0\), a nonzero term divisible by \(q^j\) forces
\(U/q^j\ge p^{r-j}\ge p^{\kappa}\). Division by \(q^j\) preserves
all the relevant coprimality conditions, because \(q\in L\),
\(q\ge p\), and \(L\cap E=\varnothing\). The upper estimate in
Lemma~\ref{an:exceptional-sieve} therefore bounds the total low-power
contribution by
\[
 (1+e_k(\kappa))UV_p\vartheta(E)\mu_L.
\]
The high-power estimate before inserting the density is unchanged:
it is at most \(U\eta p^{1-j_0}/(p-1)\). Now
\(V_p\ge1/(p-1)\) and \(\vartheta(E)\ge2^{-k}\), so this is at most
\[
 UV_p\vartheta(E)\eta2^{k+1-j_0}.
\]
Divide the sum of these bounds by \(\ell-\eta\), use
\eqref{an:conditional-condition}, and sum over carriers. This proves
the proposition, including its real-length formulation.
\end{proof}

\begin{theorem}[Exceptional coordinates with an explicit defect]
\label{an:fixed-core}
For every \(\varepsilon>0\) there is \(c_\varepsilon>0\) such that
the following holds for every integer \(k\ge0\). Suppose
\(0\le w_p\le1\) and there is a subset \(E\subseteq P\) of size at most \(k\)
such that
\[
 w_p\le\frac{c_\varepsilon}{k+1}\qquad(p\notin E).
\]
Then, uniformly in \(M\), \(I\), and the actual exceptional primes,
\begin{align}
 \Delta_{2+\varepsilon}(M,I,w)
 &\le 2^{|E|+1}
   \sum_{\substack{h\le M\\h\ E\text{-smooth}}}(Z(h)-1)_+
   \label{an:exact-core-defect}\\
 &\le2^{k+1}(1+\log_2M)^{k+1},
 \label{an:core-defect}
\end{align}
where \(Z(h)=\sum_{p\in E}w_pv_p(h)\).
\end{theorem}

\begin{proof}
Again it suffices to consider \(0<\varepsilon\le1\). We first make
the small-weight threshold quantitative. For \(u\ge16\),
\[
 e(u/2)\le3C2^{-u/4},
\]
because \(u/4\ge4\) implies
\((u/4)^{-u/8}\le2^{-u/4}\). Hence
\begin{equation}
 e_k(u)\le3C3^k2^{-u/4}+6\cdot4^k2^{-u/2}.
 \label{an:quant-sieve}
\end{equation}
Choose an integer \(R_\varepsilon\ge16\) so large that
\[
 3C2^{-R_\varepsilon/4}+6\cdot2^{-R_\varepsilon/2}
 \le\varepsilon/16.
\]
For \(u\ge16k+R_\varepsilon\), \eqref{an:quant-sieve} then gives
\(e_k(u)\le\varepsilon/16\). Choose \(c_\varepsilon>0\) satisfying
\[
 c_\varepsilon\le\varepsilon^2/128,
 \qquad
 \frac{1+\varepsilon}{c_\varepsilon}
 \ge\max(32,2R_\varepsilon+2).
\]
Set \(\eta=c_\varepsilon/(k+1)\), \(\ell=1+\varepsilon\), and
define \(r,j_0,\kappa\) as in Proposition~\ref{an:conditional-level}.
Then
\[
 r\ge32k+2R_\varepsilon,\qquad
 \kappa\ge16k+R_\varepsilon,\qquad j_0\ge k+1.
\]
In particular \(e_k(r),e_k(\kappa)\le\varepsilon/16\), and the
tail term \(\eta2^{k+1-j_0}\) is at most \(\eta\). If
\(\mu_L\le1+\varepsilon/2\), these estimates verify
\eqref{an:conditional-condition}. Explicitly, using
\(\eta\le\varepsilon/128\), the difference between its resulting
right-hand lower bound and left-hand upper bound is at least
\[
 \frac{46}{128}\varepsilon-\frac{191}{2048}\varepsilon^2>0.
\]

The case \(M=1\) is immediate. Otherwise discard all primes greater
than \(M\) before computing \(\mu_L\) or splitting into cases.
Restoring them only increases the target hinge. Removed exceptional
primes cannot occur in any \(E\)-smooth \(h\le M\), so this leaves
the smooth-part sum unchanged and can only decrease the bound in
\eqref{an:exact-core-defect}. Write
\(X_w=V+Z\), where \(V\) uses the primes outside \(E\).
If \(\mu_L\ge1+\varepsilon/2\), Lemma~\ref{an:moments} and
\eqref{an:automatic-hinge} give \(H_2\) for the full load, which is
stronger than both asserted bounds.

It remains to consider \(\mu_L\le1+\varepsilon/2\). Each positive
integer has a unique decomposition
\[
 a=ht,\qquad h\ E\text{-smooth},\qquad(t,E^{\Pi})=1.
\]
Here \(h\) is its \emph{exact} exceptional-prime part, and
\(X_w(a)=Z(h)+V(t)\). Fix \(s\ge0\). At a fixed \(h\le M\),
the prefix and target counts are over intervals of the same real length
\(M/h\). If \(s-Z(h)\ge0\), apply
Proposition~\ref{an:conditional-level} at that level; this part has
nonpositive level-count defect.

If \(s-Z(h)<0\), every target cofactor coprime to \(E^{\Pi}\) meets
the target threshold. The source count is at most the unrestricted
prefix coprimality count. Finite inclusion--exclusion shows that each
of these unrestricted counts differs from
\((M/h)\vartheta(E)\) by at most \(2^{|E|}\). Their difference is
therefore at most \(2^{|E|+1}\). Target integers with exact
exceptional part \(h>M\) are nonnegative contributions and may be
discarded. Summing gives
\begin{align}
 &\#\{a\le M:X_w(a)>s+\ell\}
   -\#\{b\in I:X_w(b)>s\} \notag\\
 &\hspace{15mm}\le
 2^{|E|+1}\#\{h\le M:h\ E\text{-smooth},\ Z(h)>s\}.
 \label{an:core-level-defect}
\end{align}
Integrating over \(s\ge1\) proves \eqref{an:exact-core-defect}.
Finally, \(Z(h)\le\Omega(h)\le\log_2M\), and there are at most
\((1+\log_2M)^{|E|}\) exceptional smooth parts \(h\le M\), by
bounding each exponent separately. This proves \eqref{an:core-defect}.
\end{proof}

\begin{corollary}[A growing-core exclusion for raw dual witnesses]
\label{an:growing-core}
Fix \(\varepsilon>0\), and consider weights in \([0,1]\). If
\(k=o(\log M/\log\log M)\) and all but
at most \(k\) weights are at most \(c_\varepsilon/(k+1)\), then
\[
 \Delta_{2+\varepsilon}(M,I,w)\le M^{o(1)}=o(\sqrt M).
\]
More specifically, fix \(\delta\in(0,1)\), let \(A\subseteq[2,M]\)
have \(n\) distinct elements, and suppose
\[
 M\le16n^2,
 \qquad k\le(1-\delta)\frac{\log n}{\log\log n}.
\]
For such weight vectors, as \(n\to\infty\),
\begin{equation}
 \sum_{a\in A}X_w(a)-\sum_{b\in I}(X_w(b)-1)_+
 \le(2+\varepsilon)n+n^{1-\delta+o(1)}.
 \label{an:dual-exclusion}
\end{equation}
\end{corollary}

\begin{proof}
The logarithm of the right side of \eqref{an:core-defect} is
\((k+1)\{\log2+\log(1+\log_2M)\}\). This is \(o(\log M)\) in
the first range. In the second range it is at most
\((1-\delta+o(1))\log n\). Finally,
\[
 \sum_{a\in A}X_w(a)
 \le(2+\varepsilon)n+
      \sum_{a=1}^{M}(X_w(a)-2-\varepsilon)_+,
\]
where distinctness is used to pass from \(A\) to the full prefix.
Combine this with Theorem~\ref{an:fixed-core}.
\end{proof}

Corollary~\ref{an:growing-core} excludes a class of raw dual vectors.
It does not assert a per-instance integer-cover bound equal to the raw
LP value plus \(o(n)\). In particular, a core of size comparable to
a positive power of \(n\) is outside its effective range.

\subsection{Dense support up to a subexponential cutoff}

\begin{theorem}[Uniform dense-support hinge inequality]
\label{an:subexp-hinge}
For every \(\varepsilon>0\) there are
\(\gamma_\varepsilon>0\) and \(M_\varepsilon\) such that
\(H_{2+\varepsilon}\) holds for all \(M\ge M_\varepsilon\), all
shifts, and all weights \(0\le w_p\le1\) supported on primes
\[
 p\le Y\le\exp(\gamma_\varepsilon\sqrt{\log M}),\qquad Y\ge2.
\]
One admissible choice is
\begin{equation}
 J_\varepsilon=\left\lceil\frac2{\eta_\varepsilon}\right\rceil,
 \qquad
 \gamma_\varepsilon=\frac1{\sqrt{16(J_\varepsilon+5)}},
 \label{an:gamma}
\end{equation}
where \(\eta_\varepsilon\) is any threshold from
Theorem~\ref{an:diffuse-theorem}. The constant
\(M_\varepsilon\) is then sufficiently large.
\end{theorem}

\begin{proof}
The empty support is trivial. Give each \(P\)-smooth positive integer
\(h\) probability \(\vartheta(P)/h\). These probabilities sum to one
by the convergent finite Euler product. Under this geometric reference
distribution,
\begin{equation}
 \mathbb E X_w=\sum_{p\in P}\frac{w_p}{p-1}\le1+\log Y.
 \label{an:model-mean}
\end{equation}
If some \(w_p\ge\eta_\varepsilon\), put \(j=\lceil2/w_p\rceil\).
Then \(j\le J_\varepsilon\), and \(X_w(h)\ge2\) whenever
\(p^j\mid h\). The difference
\((X_w-1)_+-(X_w-2)_+\) is nonnegative and equals one on this event.
Its reference mean therefore satisfies
\begin{equation}
 G:=\mathbb E\{(X_w-1)_+-(X_w-2)_+\}
 \ge p^{-j}\ge Y^{-J_\varepsilon}.
 \label{an:model-gap}
\end{equation}

We first bound the contribution of large smooth parts. Set
\[
 u=\frac{\log M}{\log Y},\qquad H=\sqrt M,\qquad
 \sigma=\frac1{4\log Y},\qquad
 R=\sum_{\substack{h>H\\h\ P\text{-smooth}}}\frac{X_w(h)}h.
\]
Since \(X_w(h)\le\log h/\log2\), Rankin's inequality gives
\begin{equation}
 R\le\frac{H^{-\sigma}}{\log2}
 \sum_{h\ P\text{-smooth}}\frac{\log h}{h^{1-\sigma}}.
 \label{an:rankin-start}
\end{equation}
For \(p\le Y\), put \(z_p=p^{-1+\sigma}\). Then
\[
 z_p\le e^{1/4}/p,
 \qquad \frac{z_p}{1-z_p}\le\frac4p,
\]
using \(e^{1/4}<4/3\). Thus
\[
 F:=\prod_{p\in P}(1-p^{-1+\sigma})^{-1}\le Y^4,
\]
because \(\log F\le4\sum_{p\le Y}1/p\le4\log Y\).
Differentiation of this finite Euler product yields
\[
 \sum_{h\ P\text{-smooth}}\frac{\log h}{h^{1-\sigma}}
 =F\sum_{p\in P}\frac{(\log p)z_p}{1-z_p}
 \le4Y^4(\log Y)^2.
\]
Here \(\sum_{p\le Y}(\log p)/p\le(\log Y)^2\), by comparison with
all integers. Since \(H^{-\sigma}=e^{-u/8}\),
\begin{equation}
 R\le\frac4{\log2}Y^4(\log Y)^2e^{-u/8}.
 \label{an:rankin-tail}
\end{equation}

Decompose each positive integer uniquely as \(ht\), where \(h\) is
its exact \(P\)-smooth part and \((t,P^{\Pi})=1\). For \(h\le H\),
both prefix and target counts at that smooth part are coprimality
counts in intervals of real length \(M/h\ge\sqrt M=Y^{u/2}\).
By Lemma~\ref{an:sieve}, their relative errors from
\((M/h)\vartheta(P)\) are at most \(e(u/2)\), provided \(u\ge8\).
The prefix contribution from \(h>H\) is at most \(MR\), since there
are at most \(M/h\) prefix multiples of \(h\). Discard the
nonnegative target contribution from these large parts.

The truncated reference gap loses at most \(R\): indeed
\(0\le(X_w-1)_+-(X_w-2)_+\le X_w\) and
\(\vartheta(P)\le1\). The sum of the two reference hinge means is
at most \(2\mathbb E X_w\). Therefore
\begin{equation}
 \frac{\Delta_2(M,I,w)}M
 \le-Y^{-J_\varepsilon}+2R+2e(u/2)(1+\log Y).
 \label{an:subexp-master}
\end{equation}

We check this comparison uniformly over the entire asserted range of
\(Y\). Write \(L=\log M\), \(y=\log Y\), and use the value
\(\gamma=\gamma_\varepsilon\) in \eqref{an:gamma}. Then
\(\log2\le y\le\gamma\sqrt L\), and the ratio of the bound in
\eqref{an:rankin-tail} to \(Y^{-J_\varepsilon}\) is
\[
 \frac4{\log2}y^2
 \exp\left((J_\varepsilon+4)y-\frac{L}{8y}\right).
\]
Both its polynomial factor and its exponent increase with \(y\).
At the upper endpoint the exponent is
\[
 \left((J_\varepsilon+4)\gamma-\frac1{8\gamma}\right)\sqrt L
 =-\gamma(J_\varepsilon+6)\sqrt L.
\]
The ratio consequently tends to zero uniformly.

For \(u\ge16\), the error function satisfies
\(e(u/2)\le3C2^{-u/4}\). The ratio of the last error in
\eqref{an:subexp-master} to \(Y^{-J_\varepsilon}\) is at most
\[
 6C(1+y)\exp\left(J_\varepsilon y-
                        \frac{(\log2)L}{4y}\right).
\]
Its exponent also increases with \(y\), and is a negative constant
times \(\sqrt L\) at the upper endpoint, since
\[
 J_\varepsilon\gamma^2<\frac1{16}<\frac{\log2}{4}.
\]
It too tends to zero uniformly. Thus, for sufficiently large \(M\),
the errors in \eqref{an:subexp-master} are dominated by the gap,
proving \(H_2\) whenever some weight is at least
\(\eta_\varepsilon\). If every weight is at most this threshold,
Theorem~\ref{an:diffuse-theorem} gives \(H_{2+\varepsilon}\) at every
length without a support restriction. The two cases prove the theorem.
\end{proof}

\subsection{The raw boxed dual and an integer-cover consequence}

We include the LP bridge and the required rounding argument to distinguish
the ensuing integer conclusion from the dual-vector exclusions above.
Let \(A\subseteq[2,M]\) be a set of \(n\) distinct sources. Write
\(\tau(A,I)\) for the least size of a subset \(B\subseteq I\) such
that \(\prod_{a\in A}a\) divides \(\prod_{b\in B}b\). Let
\(P_A\) be the set of primes occurring in \(A\). Put
\(D_p=\sum_{a\in A}v_p(a)\). The raw boxed relaxation is
\begin{equation}
 L(A,I)=\min\left\{\sum_{b\in I}x_b:
 0\le x_b\le1,\quad
 \sum_{b\in I}v_p(b)x_b\ge D_p\quad(p\in P_A)\right\}.
 \label{an:raw-lp}
\end{equation}

\begin{lemma}[Raw dual clipping and row-count rounding]
\label{an:lp-bridge}
The program \eqref{an:raw-lp} is feasible, and its value equals
\begin{equation}
 \max_{0\le w_p\le1}
 \left\{\sum_{a\in A}X_w(a)
          -\sum_{b\in I}(X_w(b)-1)_+\right\},
 \label{an:raw-dual}
\end{equation}
where the weights are on \(P_A\). If \(H_c\) holds for all these
weight vectors, then \(L(A,I)\le cn\). Furthermore,
\begin{equation}
 \tau(A,I)\le L(A,I)+|P_A|.
 \label{an:row-rounding}
\end{equation}
\end{lemma}

\begin{proof}
Distinctness gives \(\prod_{a\in A}a\mid M!\), while
\(M!\mid\prod_{b\in I}b\). Thus selecting every interval position is
feasible. Equivalently, \(\sum_{b\in I}v_p(b)\ge D_p\) for every
source prime. Finite LP duality gives the maximum in
\eqref{an:raw-dual} initially over all \(w_p\ge0\).

If some \(w_p>1\), lower it to one while keeping the other weights
fixed. For every column with \(v_p(b)>0\), its hinge is linear
throughout this change, because the final load is still at least one.
The change in the dual objective is
\[
 (w_p-1)\left(\sum_{b\in I}v_p(b)-D_p\right)\ge0.
\]
Columns with \(v_p(b)=0\) are unchanged. Clipping coordinates in turn
therefore proves \eqref{an:raw-dual}.

If \(H_c\) holds, then every vector in \eqref{an:raw-dual} has
objective at most \(cn\), since
\[
 \sum_{a\in A}X_w(a)
 \le cn+\sum_{a\in A}(X_w(a)-c)_+
 \le cn+\sum_{a=1}^{M}(X_w(a)-c)_+.
\]
For rounding, choose an optimal vertex of the feasible boxed polytope.
It has at most \(|P_A|\) coordinates strictly between zero and one.
Otherwise a nonzero perturbation supported on the fractional
coordinates could annihilate all prime rows, and sufficiently small
perturbations in both signs would preserve feasibility and the box,
contradicting extremality. Round each fractional coordinate upward to
one. The nonnegative valuation rows remain satisfied, and the cost
increases by at most \(|P_A|\), proving \eqref{an:row-rounding}.
\end{proof}

\begin{corollary}[A coefficient-two cover for subexponentially smooth sources]
\label{an:smooth-cover}
Let \(Y(M)\ge2\) be a fixed function satisfying
\[
 \log Y(M)=o(\sqrt{\log M}).
\]
Let \(A\) range over sets of \(n\) distinct integers at least two,
with \(M=\max A\), such that every source prime is at most \(Y(M)\).
Then, as \(n\to\infty\), every positive \(M\)-term interval admits
a product cover with
\[
 \tau(A,I)\le(2+o(1))n,
\]
uniformly in the source set subject to this condition and in the interval
shift. In particular this applies to
\(Y(M)=\exp((\log M)^a)\) for every fixed \(0<a<1/2\).
\end{corollary}

\begin{proof}
Fix \(\varepsilon>0\). For all sufficiently large \(M\), the cutoff
is below \(\exp(\gamma_\varepsilon\sqrt{\log M})\).
Theorem~\ref{an:subexp-hinge} and Lemma~\ref{an:lp-bridge} give
\[
 \tau(A,I)\le(2+\varepsilon)n+\pi(Y(M)).
\]
If \(M<16n^2\), then \(n>\sqrt M/4\), whereas
\(\pi(Y(M))\le Y(M)=M^{o(1)}=o(\sqrt M)\). The rounding term is
therefore \(o(n)\) in this range.

If \(M\ge16n^2\), the quadratic coefficient-two theorem of
Wang~\cite{Wang} gives the stronger exact estimate
\(\tau(A,I)\le2n\). Its hypotheses are precisely a set of \(n\)
distinct positive sources at least two, ambient length \(M=\max A\),
and an arbitrary positive interval of that length. Finally,
\(M\ge n+1\) ensures \(M\to\infty\) whenever \(n\to\infty\).
Since \(\varepsilon\) was arbitrary, the two ranges prove the
corollary.
\end{proof}

The cutoff constant \(\gamma_\varepsilon\) depends on
\(\varepsilon\). The preceding corollary therefore uses
\(\log Y=o(\sqrt{\log M})\); it does not assert coefficient two for
one fixed positive multiple of \(\sqrt{\log M}\) in the exponent at
all accuracies. The conclusion also relies on the explicit small
prime-row count in \eqref{an:row-rounding}, rather than a general
fractional-to-integer error assumption. Polynomial prime scales
\(M^\alpha\), and their interactions with a dense small-prime part,
remain outside the support theorem proved here.

\section{Exact coefficient two for separated prime bands}
\label{sec:separated-bands}

The following result permits a growing, possibly dense small-prime support.
Its source restriction is that every source outside the pure small-prime
part contains exactly one large prime occurrence. The proof assigns targets
in disjoint exact smooth-part classes and then uses one global repair phase.
It does not combine separate covers of the two source parts.

Let $A\subseteq\{2,\ldots,M\}$ be a set of $n$ distinct integers and let
$I=\{x+1,\ldots,x+M\}$, where $x\in\mathbb Z_{\ge0}$. The length $M$
is an ambient bound and need not equal $\max A$. Let $2\le Y<Z$, let $Q$
be a subset of the primes at most $Y$, and let $R$ be a finite set of primes
at least $Z$. Write
\[
 Q^\Pi=\prod_{q\in Q}q,
 \qquad \vartheta(Q)=\prod_{q\in Q}(1-1/q),
 \qquad \pi_R(T)=\#\{r\in R:r\le T\}.
\]
Both empty products equal one, and one is regarded as $Q$-smooth. Partition
the sources as $A=A_Q\sqcup A_R$, where each source in $A_Q$ is $Q$-smooth
and each source in $A_R$ has the form
\begin{equation}
 a=d(a)r(a),\qquad d(a)\ \text{$Q$-smooth},\quad r(a)\in R.
 \label{eq:sep-source-form}
\end{equation}
Thus the large prime has exponent one; sources with two large prime
occurrences are not included in this hypothesis.

We use the decreasing error function $e(\kappa)\to0$ from
Lemma~\ref{an:sieve}. In particular, for every real interval $J$ of length
$T\ge Y^\kappa$, uniformly over its shift and the choice of $Q$,
\begin{equation}
 \#\{m\in J:(m,Q^\Pi)=1\}
 \ge (1-e(\kappa))\vartheta(Q)T,
 \qquad \kappa\ge4.
 \label{eq:sep-sieve}
\end{equation}
This is the fundamental-lemma input of \cite{Lagarias,GranvilleSound}; it
includes nonintegral interval lengths. We also use absolute constants
$c_0,C_0>0$ in the standard Mertens and Chebyshev bounds
\begin{equation}
 \prod_{p\le Y}(1-1/p)\ge\frac{c_0}{\log Y},
 \qquad \pi(T)\le\frac{C_0T}{\log T}
 \quad(Y,T\ge2).
 \label{eq:sep-prime-estimates}
\end{equation}
No numerical values for the resulting separation exponent are asserted.

\begin{proposition}[A finite support comparison]
\label{prop:sep-finite}
Suppose $\kappa\ge4$, $e(\kappa)<1$, $Z\ge Y^\kappa$, and
\begin{equation}
 (1-e(\kappa))\vartheta(Q)
 \ge \sup_{T\ge Z}\frac{\pi_R(T)}T.
 \label{eq:sep-finite-condition}
\end{equation}
Then, for every source set satisfying \eqref{eq:sep-source-form},
\begin{equation}
 \tau(A,I)\le 2|A_R|+\sum_{a\in A_Q}\Omega(a).
 \label{eq:sep-cover-bound}
\end{equation}
For finite $R=\{r_1<\cdots<r_s\}\ne\varnothing$, the supremum in
\eqref{eq:sep-finite-condition} is $\max_{1\le j\le s}j/r_j$; for empty
$R$ it is zero.
\end{proposition}

\begin{proof}
For each distinct $d$ occurring in \eqref{eq:sep-source-form}, define
\[
 m_d=\#\{a\in A_R:d(a)=d\},\qquad T_d=M/d,
 \qquad B_d=\{b\in I:d\mid b,\ (b/d,Q^\Pi)=1\}.
\]
Every $b\in B_d$ has exact $Q$-part $d$: the product of all its prime
powers at primes in $Q$ equals $d$. Hence the sets $B_d$ are pairwise
disjoint, including the case $d=1$. Distinct sources with this cofactor
have distinct primes $r\le T_d$, so
\begin{equation}
 m_d\le\pi_R(T_d),\qquad T_d\ge Z.
 \label{eq:sep-fibre}
\end{equation}
The multiples of $d$ in $I$, after division by $d$, are the integer points
of $(x/d,(x+M)/d]$, which has real length $T_d$. Therefore
\eqref{eq:sep-sieve} and \eqref{eq:sep-finite-condition} give
\[
 |B_d|\ge(1-e(\kappa))\vartheta(Q)T_d
 \ge\pi_R(T_d)\ge m_d.
\]
Choose $m_d$ distinct targets from each $B_d$ and assign them to the
corresponding sources. Since the bins are disjoint, all $|A_R|$ assigned
targets are distinct. Their product is divisible by
$\prod_{a\in A_R}d(a)$.

The sum of the remaining positive prime-valuation deficits relative to
$\prod A$ is at most
\[
 |A_R|+\sum_{a\in A_Q}\Omega(a).
\]
There is one possible unit for each removed prime $r(a)$, together with
all units of the pure $Q$ sources; extra factors already supplied by the
chosen targets only help. Since $\prod A\mid M!\mid\prod I$, each still
missing unit can be supplied by an unused target. Such a repair decreases
the total positive deficit by at least one. Beginning with $|A_R|$ assigned
targets, this one global repair phase proves \eqref{eq:sep-cover-bound}.

Finally, $\pi_R(T)/T$ decreases between its jumps at the primes of $R$,
which gives the displayed finite expression for its supremum.
\end{proof}

For a particular source set, the proof only needs
$(1-e(\kappa))\vartheta(Q)T_d\ge m_d$ for its actual fibres. Condition
\eqref{eq:sep-finite-condition} is a support-only sufficient comparison.

\begin{theorem}[Absolute power separation]
\label{thm:sep-power}
There is an absolute constant $\kappa_0\ge4$ such that whenever
$Y\ge2$ and $Z\ge Y^{\kappa_0}$, every source set of the form
\eqref{eq:sep-source-form} satisfies \eqref{eq:sep-cover-bound}, for every
ambient $M\ge\max A$ and every integer shift $x\ge0$. In particular, if
$A_Q=\varnothing$, then $\tau(A,I)\le2n$.
\end{theorem}

\begin{proof}
Choose an absolute $\kappa_0$ such that
\[
 e(\kappa_0)\le\tfrac12,
 \qquad \kappa_0\ge\frac{2C_0}{c_0}.
\]
This is possible because $e(\kappa)\to0$. For every $T\ge Z\ge
Y^{\kappa_0}$, the subset relation for $Q$ and
\eqref{eq:sep-prime-estimates} give
\[
 (1-e(\kappa_0))\vartheta(Q)
 \ge\frac{c_0}{2\log Y}
 \ge\frac{C_0}{\log T}
 \ge\frac{\pi_R(T)}T.
\]
Proposition~\ref{prop:sep-finite} applies.
\end{proof}

The constant $\kappa_0$ is not computed or claimed to be practical. The
theorem is a finite statement for all parameters meeting its separation
condition. It implies the proposed asymptotic class $Z\to\infty$,
$\log Y=o(\log Z)$; for example, $Y=M^{o(1)}$ and $Z=M^\alpha$ for any
fixed $\alpha>0$ eventually satisfy the condition. The restriction to one
large prime occurrence per source remains essential to the stated result.

\begin{corollary}[Pure small-prime sources]
\label{cor:sep-pure}
Under the hypotheses of Theorem~\ref{thm:sep-power},
\begin{equation}
 \tau(A,I)\le2n+\sum_{a\in A_Q}(\Omega(a)-2)_+.
 \label{eq:sep-excess}
\end{equation}
Thus $\tau(A,I)\le2n$ if every source in $A_Q$ has at most two prime
factors with multiplicity. More generally, it is enough that
$\sum_{a\in A_Q}\Omega(a)\le2|A_Q|$. Along any sequence with
$n\to\infty$ and $\sum_{a\in A_Q}(\Omega(a)-2)_+=o(n)$, the upper
bound is $(2+o(1))n$, uniformly in the shift.
\end{corollary}

\begin{proof}
Subtract $2n=2|A_R|+2|A_Q|$ from \eqref{eq:sep-cover-bound}, and bound
each possibly negative term $\Omega(a)-2$ by its positive part.
\end{proof}

This is an upper-bound result; no matching lower coefficient is asserted.
An unrestricted pure $Q$ component may have large excess in
\eqref{eq:sep-excess}, and sources with multiple $R$ occurrences remain
outside the theorem. Consequently the argument does not settle the
unrestricted exact or asymptotic problem.

\section{Exact certificates and limitations of local reductions}
\label{root:certificates}

The companion code uses integer arithmetic and the Python standard
library. Its certificates are finite statements with explicit
hypotheses. They do not convert an unsuccessful search into a
counterexample or a universal theorem.

\subsection{A unified sufficient recognizer}

The program \texttt{exact2n\_recognition.py} keeps the ambient
length \(M\) fixed. At a current set of \(m\) sources it recognizes
the full-interval base, Wang's external quadratic base, or
Theorem~\ref{root:hall-theorem}. Otherwise it may delete a source by
one of the following proved extension rules:
\[
 \Omega(a)\le2;
\]
\[
 pq\mid a,\quad pq\ge2m-1,\quad
 D_p(A)\le q\,v_p(a/q);
\]
\[
 a=p^e q,\quad p,q\text{ prime},\quad mp^e\le M.
\]
Each valid deletion is restored with at most two additional
positions, so a successful trace proves the exact \(2|A|\) bound
for every shift. The certificate verifier starts with the original
source list and recomputes demands, cofactors, numerical witnesses,
and all Hall rows. It does not rerun the recognizer's choice heuristic;
it does share the exact arithmetic helpers. If the rules stop, the
output is \texttt{UNKNOWN}. No cover is explicitly constructed by
this recognizer.

For a concrete combination example, let \(\mathcal R\) be all 35
cubic monomials, with repetitions, on
\(\{11,13,17,19,23\}\), and set
\[
 A=\mathcal R\cup\{4096,11875\},\qquad
 n=37,\quad M=12167.
\]
Every source has at least three prime factors, counting multiplicity.
The direct Hall certificate fails: at \(d=121\) its lower bound is
\(-7\), while the fibre has size five. Nevertheless,
\[
 11875=5^4\cdot19,\qquad
 5\cdot19=95\ge73,\qquad D_5(A)=4\le19\cdot4.
\]
The two-prime rule deletes this source. All Hall rows of the remaining
36 sources then pass; their minimum margin is zero. Hence every
length-12167 interval has a cover of size at most 74 for this set.
Wang's quadratic base is inapplicable both before and after deletion,
and the prime-power-bank rule alone cannot remove this source because
\(37\cdot625>12167\).

The saved examples also include deleting the maximum while retaining
the original ambient length, two successive deletion steps, direct
base cases, and an unresolved 64-source smooth input. Their validation
checks the combined proof mechanism rather than sampling interval
shifts.

\subsection{Necessary arithmetic cores}

The program \texttt{necessary\_core.py} applies necessary conditions
for a \emph{minimal-cardinality} counterexample. It computes exact
minimum forced-union cardinalities over a period and rigorous lower
bounds for their valuations. The latter are sums of individual level
minima; they are not asserted to equal the minimum total valuation at
one common shift. It then removes candidates whose required demand
exceeds the largest possible demand of a source set of the stated
cardinality.

For example, at
\[
 n=100,\qquad M=146969=47\cdot53\cdot59,
\]
the forced 47-valuation of the maximum source is at least 112, while
the largest possible 47-demand of 100 sources from the candidate pool,
including the maximum, is at most 90. The accompanying exact trace
therefore excludes these \emph{minimal-counterexample parameters}.
It does not establish \(g(100)\le200\).

There is also an exact source-only family showing the limits of
these filters. Its certificate has
\[
 n=16384,\qquad M=n^2=268435456,
\]
and uses all 30 primes at most 113. Its source and cofactor prime
supports coincide, its common gcd is one, and its prime-support
hypergraph is connected. Every source has at least three prime
occurrences, and every prime-deletion cofactor exceeds \(M/n\).
All prime-pair products are below \(2n-1\), so the two-prime demand
tests are vacuous. The forced-union cardinality bound is at most 676,
well below the available 32766 positions, while the cofactor
reciprocal mass exceeds one and the finite Hall certificate fails.

This is a witness to the insufficiency of the listed arithmetic
conditions. It is not a counterexample to the covering conjecture:
in particular at shift zero the original sources cover themselves.
No collection of deletion covers satisfying every universal
minimality requirement is asserted to exist.

\subsection{Why unrestricted cofactor matching cannot be the missing step}

Even for sources with at least three prime factors, the particular
assignment \(a\mapsto a/P^+(a)\) need not admit a matching. An exact
example retained from the preceding work is
\[
 M=440,\quad I=\{4401,\ldots,4840\},\quad
 A=\{2k:k\le220\text{ composite},\ k\ne169\}.
\]
There are 171 sources. The union of their cofactor neighborhoods has
only 170 positions: exactly the interval integers divisible by at
least one of \(4,6,10,14,22\). Thus Hall fails for this proposed
assignment, although the original product-cover problem allows
other constructions and a much larger budget. This obstruction
concerns the proposed matching, not the \(2n\) conjecture.

\section{The remaining mathematical target}
\label{root:remaining}

The exact and asymptotic goals must be kept distinct. The structural
core defined by source demands is not automatically the same set as
the large coordinates of an optimizing dual vector. Nor can separate
covers for two source parts be combined by adding their valuation
contributions: an overlapping target is only available once.

\subsection{An exact criterion for asymptotic closure}

The raw boxed dual in Lemma~\ref{an:lp-bridge}, together with Wang's
uniform raw fractional transfer theorem \cite{Wang}, gives the
following useful precise target, also recorded in \cite{Earlier}.

\begin{proposition}[Uniform additive-defect criterion]
\label{root:closure-criterion}
The unrestricted assertion \(g(n)\le(2+o(1))n\) is equivalent to
\begin{equation}
 \forall\varepsilon>0\ \exists C_\varepsilon<\infty\
 \forall M,I,w:\qquad
 \Delta_{2+\varepsilon}(M,I,w)\le C_\varepsilon,
 \label{root:closure-target}
\end{equation}
where \(w\) has finite support and \(0\le w_p\le1\), and
\(C_\varepsilon\) depends only on \(\varepsilon\).
\end{proposition}

\begin{proof}
The asymptotic assertion implies, for each \(\varepsilon>0\), a
bound \(g(n)\le(2+\varepsilon)n+C_\varepsilon\) for all \(n\),
absorbing finitely many small cardinalities into the constant. It
bounds every boxed dual objective for every
\(A\subseteq[2,M]\). If \(\max A<M\), use a contained interval
of length \(\max A\), whose available target hinge sum is at most
that on \(I\). Maximize over all such subsets and use
\[
 \max_A\left\{\sum_{a\in A}X_w(a)-(2+\varepsilon)|A|\right\}
 =\sum_{a\le M}(X_w(a)-2-\varepsilon)_+.
\]
This proves \eqref{root:closure-target}.

Conversely, \eqref{root:closure-target} and the boxed dual give
\(L(A,I)\le(2+\varepsilon)n+C_\varepsilon\) uniformly. For
completeness, the specific external transfer inequality being used
is as follows. Put \(\ell(0)=0\), and let \(\ell(m)\) be the
supremum of the raw boxed values on original \(m\)-source
instances. Then
\[
 K_n=\left\lceil\max_{0\le m\le n}
        \{\ell(m)+2(n-m)\}\right\rceil,
 \qquad g(n)\le K_n+\pi(2K_n)-1.
\]
Since \(\ell(m)\le(2+\varepsilon)m+C_\varepsilon\), this yields
\(g(n)\le(2+\varepsilon)n+C_\varepsilon+o(n)\).
Let \(\varepsilon\) tend to zero.
\end{proof}

\subsection{Rounding is already small on the reduced obstruction class}

For the minimal asymptotic-obstruction sequences constructed above,
all source primes are below \(K<10n\). Hence
\[
 |P_A|\le\pi(10n)=O(n/\log n),\qquad
 \tau(A,I)\le L(A,I)+O(n/\log n),
\]
by Lemma~\ref{an:lp-bridge}. This conclusion concerns these reduced
instances; it does not assert small rounding loss for an arbitrary
original instance.

Consequently a more focused sufficient next theorem is a uniform
fractional bound
\[
 L(A,I)\le(2+\delta)n+o(n)
\]
on this explicit structural class, for every fixed \(\delta>0\).
The error must be uniform in the source set, shift, and additive
minimality parameter. Indeed, use the minimality budget
\((2+2\delta)n+C\), with \(2+2\delta<10\). As \(C\to\infty\),
the source cardinality tends to infinity, and the displayed estimates
would give \(\tau(A,I)<(2+2\delta)n\), a contradiction. Thus the
outstanding issue on this class is its joint fractional valuation
bound. That bound is not proved here.

\subsection{A concrete order of attack}

The following are mathematical milestones, rather than a claim
that the remaining work is routine.

\begin{enumerate}
\item \textbf{Range and proof audit: completed.} The submitted
\(n^{5/2}\) proof is reconstructed, its essential restriction is
identified, and the independently checked quadratic input sets the
remaining range to \eqref{root:range}.
\item \textbf{Arbitrarily many small weights: completed.} The
diffuse theorem controls this class at every scale, including one
arbitrary exceptional coordinate. No bound on the number of prime
factors is assumed.
\item \textbf{A controlled large-weight block: completed within
its stated range.} The explicit defect allows a fixed or sufficiently
slowly growing number of exceptions. It does not handle a block as
large as the source-demand core in the combinatorial reduction.
\item \textbf{Dense small-prime sources and separated bands:
completed within their stated classes.} The smooth-source theorem
and the exact separated-band theorem are actual covering results.
Neither covers arbitrary combinations of adjacent prime scales.
\item \textbf{Growing mixed cores: open.} The next substantive
target is a uniform certificate handling substantial weights on
intermediate prime scales while retaining a growing dense small-prime
part. A uniform bound \eqref{root:closure-target} would close the
asymptotic question. An exact \(2n\) construction would additionally
have to control the integer position budget without an asymptotic
rounding allowance.
\end{enumerate}

One quantitative boundary of the present proof is explicit. In the
smooth-support argument, the sieve and tail errors behave like
\(\exp(-c\log M/\log Y)\), whereas the guaranteed model gap is
\(Y^{-J_\varepsilon}\). Their comparison works when
\((\log Y)^2\ll\log M\). It does not work at a general polynomial
cutoff \(Y=M^\alpha\). The separated-band argument handles one
part of the interaction by producing disjoint exact-part bins, but
it requires a large gap between the two prime ranges and one large
occurrence per source.

The package therefore contains verified partial theorems and a
precise remaining target, not a proof that problem 708 is solved or
that a numerical fraction of it has been completed.


\clearpage
\begin{thebibliography}{9}
\bibitem{Chen}
H. Chen, \emph{Explicit upper bounds for the Erd\H{o}s--Sur\'anyi
function \(g(n)\)}, public working manuscript and accompanying
certificates. \href{https://github.com/chy4pro/erdos-708-explicit-upper-bounds}{Primary repository}.
Consulted 23 November 2026.

\bibitem{Wang}
Z. Wang, \emph{Fractional-to-Integer Transference and Five-Prime
Bounds for Interval Divisor Covers}, public working manuscript,
version 1.0.0, 21 October 2026. In particular Appendix A,
``The quadratic coefficient-two gate.''
\href{https://github.com/science-prince/erdos-708-interval-divisor-covers}{Primary repository}.
Consulted 23 November 2026.

\bibitem{Lagarias}
J. C. Lagarias, \emph{Math 675: Combinatorial Sieve: Fundamental
Lemma}, theorem and interval example, pp. 1--2.
\href{https://websites.umich.edu/lagarias/675notes/Estimate5-2015.pdf}{University of Michigan lecture notes}.

\bibitem{GranvilleSound}
A. Granville and K. Soundararajan, \emph{Multiplicative number
theory: The pretentious approach}, Chapter 1, Theorem 1.4.2.
\href{https://dms.umontreal.ca/~andrew/PDF/Book.To1.4.pdf}{Author-hosted text}.

\bibitem{TaoMertens}
T. Tao, \emph{Mertens' theorems}, 11 December 2013.
\href{https://terrytao.wordpress.com/2013/12/11/mertens-theorems/}{Author's exposition}.

\bibitem{Earlier}
\emph{Partial Results on Erd\H{o}s Problem 708: Divisor
certificates, exact covers for rough source sets, and a precise
remaining obstruction}, earlier AI-assisted working note prepared
for Denis Paliy, 23 November 2026, 14 pp. Included as
\texttt{prior\_results/erdos708\_partial\_results.pdf} in this package.
Internal research document; no publication or priority claim.
\end{thebibliography}
\end{document}
